\section{相关工作与差异}

围绕区域或场群风电预测中的时间变化和场站间空间相关性，已有研究从风速序列建模、规则网格嵌入和图结构建模等方向展开。早期深度学习方法利用 CNN-LSTM 建模风速时序特征 \cite{ref1}，也有工作将风机嵌入规则网格，使卷积网络能够处理空间分布 \cite{ref2}。面向风电场群，Wang 等在训练中自适应调整图拓扑，并结合时空图卷积、改进多头注意力和多任务输出 \cite{ref3}；Li 等将场站划分为层级子区域，并将图建模与面向极值的校正机制结合 \cite{ref4}；Liu 等识别功率波动模式，并为不同模式训练相应的时空图预测器 \cite{ref5}。Yang 等构造多种全局信息图，根据动态相关状态切换自适应图卷积模型 \cite{ref6}；后续工作通过节点嵌入动态更新场站空间相关性 \cite{ref7}，或利用图偏差注意力学习图结构，并针对负荷峰谷时段实施动态误差校正 \cite{ref8}。另有研究构造波动相关图，并依据 NWP 有效趋势切换图结构，以同时支持确定性和概率预测 \cite{ref9}；D-stacking 将时空图特征与双重多模型融合结合，用于滚动多步预测 \cite{ref10}；趋势分类和空间信息集成被用于组织日前场群预测输入 \cite{ref11}；全局感知动态时空图持续学习框架则用于适应变化的图输入并避免遗忘历史动态 \cite{ref12}。这些研究说明，多场站功率动态和空间耦合具有重要价值，但其重点主要是历史功率表示或场群图结构构造，预测起点真实功率状态应如何约束未来气象读取仍缺少显式刻画。

围绕多步风电预测中的未来气象驱动，M2TNet 通过多任务 Transformer 联合学习多模态表示和辅助任务 \cite{ref13}，M2STAN 则利用时空注意力将多模态融合扩展到多场站预测 \cite{ref14}。针对 NWP 信息的局限性，已有工作使用序列级注意力处理气象变量的时间滞后 \cite{ref15}，并针对 NWP 风速偏移设计误差容忍策略 \cite{ref16}；深度贝叶斯模型用于量化预测不确定性 \cite{ref_liu_bayesian_uncertainty}。另有研究融合多源时空特征进行场群预测 \cite{ref17}，或将空间相关性与校正后的 NWP 结合 \cite{ref18}。对于网格化气象场，3D CNN 可直接从 NWP 数据立方体中提取特征 \cite{ref19}；WindDragon 提供了自动化区域深度学习框架 \cite{ref20}；深度 CNN 编码还被用于结合保形回归森林开展区域概率预测 \cite{ref21}。本文采用 Pangu-Weather 作为未来气象来源 \cite{ref22}，并将其视为预测日内功率变化的主驱动信息。

Informer 为长序列时间序列预测引入高效注意力机制 \cite{ref23}。PatchTST 将时间序列表示为时间片段 \cite{ref24}，Crossformer 则建模多变量序列中的跨维依赖 \cite{ref25}；Non-stationary Transformers 显式处理时间序列的分布非平稳性 \cite{ref26}。在风电预测任务中，特征-时间 Transformer 用于实现可解释的多变量建模 \cite{ref_liu_feature_temporal}，TDFormer 则通过多重依赖解耦提升多步预测精度 \cite{ref_liu_tdformer}。这些方法说明，注意力机制适合处理多步预测中的长时依赖；业务化风电预测中，未来气象驱动与预测起点运行状态之间的信息分工仍需进一步显式刻画。

已有多源融合研究通常侧重气象、功率和空间信息之间的特征互补，较少关注预测起点功率状态对未来气象读取的约束作用。本文从预测起点前的场站级实测功率中生成动态上下文表示，并将其引入气象记忆读取，使各预测步的气象信息提取受到当前运行状态的约束。
